% Assignment source: ne630/homework/html/HW07.html.
% Legacy foundation: ne630_problems/lesson_07/statement.tex @ e1fba66.
% Problem 1 retains the legacy solution while updating part (e) to the OpenMC
% 294 K values requested in Fall 2026. The missing square in the legacy
% definition of alpha is corrected; the worked examples already used it.

% --------------------------------------------------------------------------
% Problem 1
% --------------------------------------------------------------------------
\noindent{\bf Problem 1}. Neutrons scatter elastically at $1.0~\mathrm{MeV}$.
After one scattering collision, determine the fraction of the neutrons that
will have energies less than $0.5~\mathrm{MeV}$ if they scatter from each of
the following:
\begin{enumerate}[label=(\alph*)]
\item Hydrogen
\item Deuterium
\item Carbon-12
\item Uranium-238
\item Light water
\end{enumerate}
For part (e), obtain the H-1 and O-16 elastic-scattering cross sections at
$1.0~\mathrm{MeV}$ using OpenMC's Python nuclear-data interface
(\texttt{openmc.data.IncidentNeutron}). Use the elastic-scattering reaction
(ENDF MT 2) and the \texttt{"294K"} cross section.

\begin{adaptedsolution}
The probability that a neutron of energy $E$ will scatter down to an energy
between $E'$ and $E'+dE'$ after colliding with a target nucleus of mass number
$A$ is given by
\[
 p(E\to E')\,dE'=
 \begin{cases}
 \dfrac{1}{(1-\alpha)E}\,dE',&\alpha E\leq E'\leq E,\\
 0,&\text{otherwise},
 \end{cases}
 \tag{FNRP 2.47}
\]
where
\[
 \alpha=\left(\frac{A-1}{A+1}\right)^2.
\]
Then, the probability that a neutron of energy $E$ scatters to an energy
equal to or less than $E_0$ is
\[
 P(E'\leq E_0)=
 \begin{cases}
 \displaystyle\int_{\alpha E}^{E_0}
 \frac{1}{(1-\alpha)E}\,dE'
 =\dfrac{E_0-\alpha E}{(1-\alpha)E},&E_0>\alpha E,\\[3pt]
 0,&\text{otherwise}.
 \end{cases}
\]
Thus, for $E=1~\mathrm{MeV}$ and $E_0=0.5~\mathrm{MeV}$, we have for each
target nucleus
\begin{enumerate}[label=(\alph*)]
\item $\alpha_{\mathrm H}=(1-1)^2/(1+1)^2=0$, and
\[
 \boxed{P(E'\leq0.5~\mathrm{MeV})=\frac{0.5}{1}=0.5}.
\]

\item $\alpha_{\mathrm D}=(2-1)^2/(2+1)^2=1/9$, and
\[
 \boxed{P(E'\leq0.5~\mathrm{MeV})
 =\frac{0.5-(1/9)1}{(1-1/9)1}=0.4375}.
\]

\item $\alpha_{\mathrm C}=(12-1)^2/(12+1)^2\approx0.716$ and, hence,
$0.5=E_0<\alpha_{\mathrm C}E=0.716$, so
\[
 \boxed{P(E'\leq0.5~\mathrm{MeV})=0}.
\]

\item $\alpha_{\mathrm U}=(238-1)^2/(238+1)^2\approx0.983$, so again
\[
 \boxed{P(E'\leq0.5~\mathrm{MeV})=0}.
\]
\end{enumerate}

For (e), we have a mixture of two nuclides and, hence, we need
\[
 p(E\to E')=\frac{1}{\Sigma_s(E)}
 \sum_i\Sigma_{s i}(E)p_i(E\to E').
 \tag{FNRP 2.53}
\]
For $1~\mathrm{MeV}$ neutrons, the 294~K OpenMC elastic-scattering cross
sections are
\[
 \sigma_s^{\isoA{H}{1}}=4.249507~\mathrm b,
 \qquad
 \sigma_s^{\isoA{O}{16}}=8.125790~\mathrm b.
\]
Because $\alpha_{\isoA{O}{16}}=(15/17)^2\approx0.779$, a $1~\mathrm{MeV}$
neutron cannot scatter off \isoA{O}{16} to energies below
$0.5~\mathrm{MeV}$. We already found that half of $1~\mathrm{MeV}$ neutrons
scatter off \isoA{H}{1} to energies below $0.5~\mathrm{MeV}$. Hence, we just
need
\[
 \frac{\Sigma_s^{\isoA{H}{1}}}{\Sigma_s^{\text{water}}}
 =\frac{2n_w\sigma_s^{\isoA{H}{1}}}
 {2n_w\sigma_s^{\isoA{H}{1}}+n_w\sigma_s^{\isoA{O}{16}}}
 =\frac{2(4.249507)}{2(4.249507)+8.125790}
 \approx0.511225,
\]
where $n_w$ is the number density of water molecules. Therefore,
\[
 \boxed{P(E'\leq0.5~\mathrm{MeV})
 =0.511225(0.5)\approx0.255612}.
\]
\end{adaptedsolution}

\newpage

% --------------------------------------------------------------------------
% Problem 2: solution language retained from the legacy source. One omitted
% line of algebra is displayed explicitly in place of the source's ellipsis.
% --------------------------------------------------------------------------
\noindent{\bf Problem 2}. Show that, for elastic scattering,
\[
 \overline{E-E'}\equiv
 \int_0^E(E-E')p(E\to E')\,dE'
 =\frac{(1-\alpha)E}{2}.
\]

\vspace{0.5cm}
\ifsolved
{
\color{purple}

\noindent{\bf SOLUTION}:

Given a probability density function $p(x)$, recall that the expected value
of some function $f(x)$ is given by $\bar{f}=\int f(x)p(x)\,dx$, where the
integral is over the domain of $x$. Here, our density is $p(E\to E')$, which
can also be written as $p(E';E)$ (i.e., the probability for $E'$ given $E$ as
a known parameter). Then, with $f(E')=E-E'$, we have
\begin{align*}
 \overline{E-E'}
 &=\int_0^E(E-E')p(E\to E')\,dE'\\
 &=\int_{\alpha E}^E(E-E')\frac{1}{(1-\alpha)E}\,dE'\\
 &=\left[\frac{E'}{1-\alpha}
 -\frac{E'^2}{2(1-\alpha)E}\right]_{\alpha E}^E\\
 &=\frac{E}{1-\alpha}
 \left(\frac12-\alpha+\frac{\alpha^2}{2}\right)\\
 &=\boxed{\frac{(1-\alpha)E}{2}}.
\end{align*}

}
\else
% NOTHING
\fi

\clearpage
